% =====================================================================
%  第 4 章 希尔伯特空间上的算子 —— ch04b（§4.3–§4.6）
%  页段：PDF 223–241（原书 p.222–240）
% =====================================================================

% [原文 p.222]
\section{斯通定理}

量子力学中，量子态的动力学由薛定谔方程
\[
\ii\hbar\frac{\pa}{\pa t}\psi=H\psi
\]
描述。假设哈密顿算子 $H$ 不依赖 $t$，则该方程形式地解为
\[
\psi(t)=\ee^{-\ii tH/\hbar}\psi_0,
\]
其中 $\psi_0=\psi(t=0)$ 是 $t=0$ 时刻的初始态。本小节讨论量子演化 $\ee^{-\ii tH/\hbar}$。基本结果是斯通定理，它把量子动力学与自伴算子联系起来。

\subsection{单参数酉群}

\begin{definition}{4.3.1}
  设 $\mathcal{H}$ 是希尔伯特空间。$\mathcal{H}$ 上的\emph{单参数酉群}是指一族酉算子 $\{U(t)\}_{t\in\R}$，满足
  \[
  \text{①}\quad U(0)=I,
  \]
  \[
  \text{②}\quad U(s)U(t)=U(s+t),\qquad \forall s,t\in\R.
  \]
  若对任意 $\varphi\in\mathcal{H}$ 与 $t_0\in\R$ 都有
  \[
  \lim_{t\to t_0}\norm{U(t)\varphi-U(t_0)\varphi}=0,
  \]
  则称它是\emph{强连续的}。
\end{definition}

\begin{definition}{4.3.2}
  设 $U(t)$ 是 $\mathcal{H}$ 上的强连续单参数酉群。定义算子
  \[
  A\varphi=\lim_{t\to 0}\frac{1}{\ii}\frac{U(t)\varphi-\varphi}{t},
  \]
  其定义域为 $\Dom(A)=\{\varphi\in\mathcal{H}\mid A\varphi\ \text{存在}\}$。称 $A$ 为 $U(t)$ 的\emph{无穷小生成元}。
\end{definition}

有限维情形下，$U(t)$ 是酉矩阵的单参数群，此时总有
\[
U(t)=\ee^{\ii tA},
\]
其中无穷小生成元 $A$ 是厄米矩阵。我们的目标是把这一结果推广到希尔伯特空间 $\mathcal{H}$ 上的无穷维情形。

\begin{definition}{4.3.3}
  设 $A$ 是 $\mathcal{H}$ 上的自伴算子。给定 $t\in\R$，定义
  \[
  \ee^{\ii tA}=\int_\R \ee^{\ii t\lambda}\,\rd\mathbb{P}_A(\lambda),
  \]
  其中对有界函数 $f_t(\lambda)=\ee^{\ii t\lambda}$ 作博雷尔函数演算。
\end{definition}

% [原文 p.223]
\subsection{斯通定理}

\begin{theorem}{4.3.4}（斯通定理）
  强连续单参数酉群与自伴算子之间存在一一对应：
  \[
  \{\text{强连续单参数酉群 }U(t)\}\ \longleftrightarrow\ \{\text{自伴算子 }A\}.
  \]
  $A$ 是 $U(t)$ 的无穷小生成元，且 $U(t)=\ee^{\ii tA}$。
\end{theorem}

\begin{proof}
  （$\Leftarrow$）设 $A$ 是自伴算子。则算子
  \[
  U(t)=\ee^{\ii tA}=\int_\R \ee^{\ii t\lambda}\,\rd\mathbb{P}_A(\lambda)
  \]
  构成 $\mathcal{H}$ 上的单参数酉群。由命题 4.2.25，对任意 $\varphi\in\mathcal{H}$，
  \begin{align*}
  \norm{U(t)\varphi-U(s)\varphi}^2
  &=\langle\varphi|(U(t)^*-U(s)^*)(U(t)-U(s))\varphi\rangle\\
  &=\int_\R\abs{\ee^{\ii t\lambda}-\ee^{\ii s\lambda}}^2\,\rd\mu_{A,\varphi}(\lambda).
  \end{align*}
  由于 $\abs{\ee^{\ii t\lambda}-\ee^{\ii s\lambda}}^2$ 有界、$\mu_{A,\varphi}(\lambda)$ 是有限测度，由控制收敛定理
  \[
  \lim_{t\to s}\norm{U(t)\varphi-U(s)\varphi}^2=0.
  \]
  这证明 $U(t)$ 强连续。

  设 $\widetilde{A}$ 是 $U(t)$ 的无穷小生成元。我们验证 $\widetilde{A}=A$。对任意 $\varphi\in\Dom(A)$，
  \[
  \left\|\frac{U(t)\varphi-\varphi}{\ii t}-A\varphi\right\|^2
  =\int_\R\left|\frac{\ee^{\ii\lambda t}-1}{\ii t}-\lambda\right|^2\,\rd\mu_{A,\varphi}.
  \]
  由于 $\varphi\in\Dom(A)$，
  \[
  \norm{A\varphi}^2=\int_\R\lambda^2\,\rd\mu_{A,\varphi}<\infty.
  \]
  又有 $\left|\frac{\ee^{\ii\lambda t}-1}{\ii t}\right|\le\lambda$。由控制收敛定理，
  \[
  \lim_{t\to 0}\left\|\frac{U(t)\varphi-\varphi}{\ii t}-A\varphi\right\|^2
  =\lim_{t\to 0}\int_\R\left|\frac{\ee^{\ii\lambda t}-1}{\ii t}-\lambda\right|^2\rd\mu_{A,\varphi}=0.
  \]
  这说明 $\varphi\in\Dom(\widetilde{A})$，且
  \[
  A\subseteq\widetilde{A}.
  \]
  另一方面，设 $\varphi,\psi\in\Dom(\widetilde{A})$。则
  \begin{align*}
  \langle\varphi|\widetilde{A}\psi\rangle
  &=\lim_{t\to 0}\left\langle\varphi\middle|\frac{U(t)\psi-\psi}{\ii t}\right\rangle
  =-\lim_{t\to 0}\left\langle\frac{U(-t)\varphi-\varphi}{\ii t}\middle|\psi\right\rangle\\
  &=\langle\widetilde{A}\varphi|\psi\rangle.
  \end{align*}
  于是 $\widetilde{A}$ 是对称算子，给出 $A$ 的一个对称扩张。由 $A$ 自伴即得
  \[
  \widetilde{A}=A.
  \]

  （$\Rightarrow$）反之，设 $U(t)$ 是 $\mathcal{H}$ 上的强连续单参数酉群。令
  \[
  A=\lim_{t\to 0}\frac{U(t)-I}{\ii t}
  \]
% [原文 p.224]
  为无穷小生成元。我们分几步证明 $A$ 是自伴算子。

  i) $A$ 是稠密定义的。

  给定任意紧支光滑函数 $f$，可定义有界算子
  \[
  U_f=\int_\R f(t)U(t)\,\rd t,
  \]
  它由有界半双线性形式
  \[
  \langle\varphi|U_f\psi\rangle=\int_\R f(t)\langle\varphi|U(t)\psi\rangle\,\rd t
  \]
  确定。直接验证可得，对任意 $\varphi\in\mathcal{H}$，
  \[
  \lim_{t\to 0}\frac{U(t)U_f\varphi-U_f\varphi}{t}=-U_{f'}\varphi.
  \]
  因此对任意 $f\in C_c^\infty(\R)$，
  \[
  \Ran U_f\subset\Dom(A).
  \]
  现在取 $f_n\in C_c^\infty(\R)$ 使得
  \[
  f_n\ge0,\qquad \supp f_n\subset\left[-\frac1n,\frac1n\right],\qquad \int_\R f_n=1.
  \]
  则对任意 $\varphi\in\mathcal{H}$，
  \[
  \norm{U_{f_n}\varphi-\varphi}
  =\left\|\int_\R f_n(t)(U(t)\varphi-\varphi)\,\rd t\right\|
  \le\sup_{t\in[-\frac1n,\frac1n]}\norm{U(t)\varphi-\varphi}.
  \]
  由于 $U(t)$ 强连续，
  \[
  \lim_{n\to\infty}\norm{U_{f_n}\varphi-\varphi}=0,
  \]
  故 $\varphi$ 可由 $U_{f_n}\varphi\in\Dom(A)$ 逼近。这证明 $\Dom(A)$ 稠密。

  ii) $A$ 是对称的。证明与上面（$\Leftarrow$）部分类似。

  iii) $\Dom(A)$ 在 $U(t)$ 下不变。

  对任意 $\varphi\in\Dom(A)$，
  \[
  \lim_{s\to 0}\frac{U(s)U(t)\varphi-U(t)\varphi}{\ii s}
  =\lim_{s\to 0}\frac{U(t)(U(s)\varphi-\varphi)}{\ii s}
  =U(t)A\varphi
  \ \Longrightarrow\ U(t)\varphi\in\Dom(A).
  \]
  特别地，
  \[
  -\ii\frac{\rd}{\rd t}(U(t)\varphi)=U(t)A\varphi=AU(t)\varphi.
  \]

  iv) $A$ 本质自伴，即 $\Ker(A^*-\ii)=\{0\}$。

  设 $\varphi\in\Ker(A^*-\ii)$，即 $A^*\varphi=\ii\varphi$。对任意 $\psi\in\Dom(A)$，
  \begin{align*}
  \ii\frac{\rd}{\rd t}\langle U(t)\psi|\varphi\rangle
  &=\left\langle-\ii\frac{\rd}{\rd t}U(t)\psi\middle|\varphi\right\rangle
  =\langle AU(t)\psi|\varphi\rangle\\
  &=\langle U(t)\psi|A^*\varphi\rangle
  =\ii\langle U(t)\psi|\varphi\rangle,
  \end{align*}
  从而
  \[
  \langle U(t)\psi|\varphi\rangle=c\,\ee^{t},\qquad\forall t\in\R.
  \]
  由于
  \[
  \abs{\langle U(t)\psi|\varphi\rangle}\le\norm{U(t)\psi}\norm{\varphi}=\norm{\psi}\norm{\varphi}
  \]
% [原文 p.225]
  有界，必须有 $c=0$，于是
  \[
  \langle\psi|\varphi\rangle=0,\qquad\forall\psi\in\Dom(A).
  \]
  由 $\Dom(A)$ 稠密得 $\varphi=0$。

  v) $A$ 自伴。

  设 $\overline{A}$ 是 $A$ 的闭包。由 iv) 知 $\overline{A}$ 是自伴算子。令
  \[
  \widetilde{U}(t)=\ee^{\ii t\overline{A}}.
  \]
  对 $\varphi\in\Dom(A)\subset\Dom(\overline{A})$，考虑
  \[
  \xi(t)=\widetilde{U}(t)\varphi-U(t)\varphi,
  \]
  它满足
  \[
  -\ii\frac{\rd}{\rd t}\xi(t)=A\xi(t).
  \]
  因此
  \[
  -\ii\frac{\rd}{\rd t}\langle\xi(t)|\xi(t)\rangle
  =\langle-A\xi(t)|\xi(t)\rangle+\langle\xi(t)|A\xi(t)\rangle=0.
  \]
  这说明
  \[
  \norm{\xi(t)}^2=\norm{\xi(0)}^2=0,
  \]
  于是
  \[
  \widetilde{U}(t)\varphi=U(t)\varphi,\qquad\forall\varphi\in\Dom(A).
  \]
  由此 $\widetilde{U}(t)=U(t)$ 在稠密的 $\Dom(A)$ 上成立，故
  \[
  \widetilde{U}(t)=U(t),\qquad\forall t\in\R.
  \]
  由（$\Leftarrow$）部分，$\overline{A}$ 是 $\widetilde{U}(t)$ 的无穷小生成元。于是 $\overline{A}=A$。
\end{proof}

\takeaway{本节证明斯通定理：强连续单参数酉群与自伴算子一一对应，前者是后者的指数，后者是前者的无穷小生成元。两个方向分别用谱演算加控制收敛、以及稠密性与本质自伴性推出。}

\section{斯通–冯·诺依曼定理}

\subsection{外尔关系}

$\R^{2n}$ 上的正则对易关系由 $2n$ 个算子 $\{Q^k,P_k\}_{k=1}^n$ 表示，形式为
\[
\begin{cases}
[Q^k,P_j]=\ii\hbar\delta^k_j,\\[2mm]
[Q^k,Q^j]=[P_k,P_j]=0,
\end{cases}
\qquad \forall j,k=1,\cdots,n.
\]
表示正则对易关系的标准方式是 $L^2(\R^n)$ 上的薛定谔表示：
\[
(Q^k\varphi)(\mathbf{x})=x^k\varphi(\mathbf{x}),\qquad
(P_k\varphi)(\mathbf{x})=-\ii\hbar\frac{\pa}{\pa x^k}\varphi(\mathbf{x}).
\]
这里 $Q^k$ 是自伴的乘法算子，$P_k$ 也自伴（经傅里叶变换）。注意所有 $Q^k$ 与 $P_k$ 都无界。事实证明这是一般现象。

% [原文 p.226]
为简单起见，我们考虑 $\R^2$ 上只有两个算子 $Q,P$ 的情形。在薛定谔表示中，$L^2(\R)$ 上的单参数酉群 $U(t)=\ee^{\ii tP}$ 与 $V(t)=\ee^{\ii tQ}$ 是
\begin{align*}
(U(t)\varphi)(x)&=\varphi(x+\hbar t),\\
(V(t)\varphi)(x)&=\ee^{\ii tx}\varphi(x).
\end{align*}
它们满足下述\emph{外尔关系}
\[
U(t)V(s)=\ee^{\ii\hbar st}V(s)U(t),
\]
该关系可视为海森堡群的一个表示（见 §5.3）。形式上，外尔关系是正则对易关系借助 Baker–Campbell–Hausdorff 公式写成的指数形式；但我们处理的是无界算子，若不加额外的定义域假设，两者并不等价。

\subsection{斯通–冯·诺依曼定理}

\begin{theorem}{4.4.1}（斯通–冯·诺依曼定理）
  海森堡群 $H_n$ 的任何不可约酉表示，若中心作用为 $\ee^K=\ee^{\ii/\hbar}$，都与 $L^2(\R^n)$ 上的薛定谔表示等价。
\end{theorem}
\tutornote{原文以 $e^K$ 表示海森堡群中心元素的像，$K$ 为中心元；这里的指数上没有 $\hbar$，与 §5.3 的约定一致。}

\begin{proof}
  为简单起见，设 $n=1$；证明可推广到一般的 $n$。海森堡群的相应表示由希尔伯特空间 $\mathcal{H}$ 上的单参数酉群 $U(t)$、$V(t)$ 组成，它们满足外尔关系
  \[
  U(t)V(s)=\ee^{\ii\hbar st}V(s)U(t).
  \]
  给定任意施瓦茨函数 $\varphi(t,s)\in\mathcal{S}(\R^2)$，可构造有界算子
  \[
  P_\varphi=\int_{\R^2}\varphi(t,s)U(t)V(s)\,\rd t\,\rd s.
  \]
  注意
  \begin{align*}
  P_\varphi^*
  &=\int_{\R^2}\overline{\varphi(t,s)}V(-s)U(-t)\,\rd t\,\rd s\\
  &=\int_{\R^2}\overline{\varphi(-t,-s)}V(s)U(t)\,\rd t\,\rd s\\
  &=\int_{\R^2}\overline{\varphi(-t,-s)}\ee^{-\ii\hbar st}U(t)V(s)\,\rd t\,\rd s.
  \end{align*}
  因此
  \[
  P_\varphi^*=P_\varphi
  \quad\text{当且仅当}\quad
  \varphi(-t,-s)=\ee^{-\ii\hbar st}\overline{\varphi(t,s)}.
  \]
  对两个 $\varphi_1,\varphi_2\in\mathcal{S}(\R^2)$，
  \begin{align*}
  P_{\varphi_1}\circ P_{\varphi_2}
  &=\int_{\R^4}\varphi_1(t_1,s_1)\varphi_2(t_2,s_2)U(t_1)V(s_1)U(t_2)V(s_2)\,\rd t_1\rd s_1\rd t_2\rd s_2\\
  &=\int_{\R^4}\varphi_1(t_1,s_1)\varphi_2(t_2,s_2)\ee^{-\ii\hbar s_1t_2}U(t_1)U(t_2)V(s_1)V(s_2)\,\rd t_1\rd s_1\rd t_2\rd s_2\\
  &=\int_{\R^4}\varphi_1(t_1,s_1)\varphi_2(t_2,s_2)\ee^{-\ii\hbar s_1t_2}U(t_1+t_2)V(s_1+s_2)\,\rd t_1\rd s_1\rd t_2\rd s_2\\
  &=\int_{\R^2}(\varphi_1\#\varphi_2)(t,s)U(t)V(s)\,\rd t\,\rd s,
  \end{align*}
% [原文 p.227]
  其中
  \[
  (\varphi_1\#\varphi_2)(t,s)=\int_{\R^2}\varphi_1(u,v)\varphi_2(t-u,s-v)\ee^{\ii\hbar(u-t)v}\,\rd u\,\rd v.
  \]
  现在考虑显式函数
  \[
  \boxed{\ \Phi(t,s)=\frac{\hbar}{2\pi}\ee^{-\frac{\ii}{2}\hbar ts}\ee^{-\frac{\hbar}{4}(s^2+t^2)}\ }
  \]
  对任意 $a,b\in\R$，记 $\delta_{a,b}$ 为 $\delta$ 函数 $\delta_{a,b}(t,s)=\delta(t-a)\delta(s-b)$。我们有
  \[
  (\delta_{a,b}\#\Phi)(t,s)=\Phi(t-a,s-b)\ee^{\ii\hbar(a-t)b},
  \]
  以及
  \begin{align*}
  &\bigl(\Phi\#(\delta_{a,b}\#\Phi)\bigr)(t,s)\\
  &=\int_{\R^2}\Phi(u,v)\Phi(t-u-a,s-v-b)\ee^{\ii\hbar(a+u-t)b}\ee^{\ii\hbar(u-t)v}\,\rd u\,\rd v\\
  &=\frac{\hbar^2}{(2\pi)^2}\int_{\R^2}
  \ee^{-\frac{\hbar}{4}\left[u^2+v^2+(t-u-a)^2+(s-v-b)^2\right]}\\
  &\qquad\qquad\qquad\times
  \ee^{\frac{\ii}{2}\hbar\left[(s+b)u-(t+a)v-(t-a)(s+b)\right]}\,\rd u\,\rd v\\
  &=\frac{\hbar^2\ee^{-\frac{\hbar}{4}\left[(t-a)^2+(s-b)^2\right]}}{(2\pi)^2}
  \int_{\R^2}\ee^{-\frac{\hbar}{2}\left(u-\frac{(t-a)+\ii(s+b)}{2}\right)^2}
  \ee^{-\frac{\hbar}{2}\left(v-\frac{(s-b)-\ii(t+a)}{2}\right)^2}\\
  &\qquad\qquad\qquad\times
  \ee^{-\frac{\hbar}{2}[ta+sb+\ii ts-\ii ab]}\,\rd u\,\rd v\\
  &=\frac{\hbar}{2\pi}\ee^{-\frac{\hbar}{4}(t^2+s^2+a^2+b^2)-\frac{\ii}{2}\hbar(ts-ab)}
  =\ee^{-\frac{\hbar}{4}(a^2+b^2)+\frac{\ii\hbar}{2}ab}\Phi(t,s).
  \end{align*}
  由此可得
  \[
  P_\Phi U(a)V(b)P_\Phi=\ee^{-\frac{\hbar}{4}(a^2+b^2)+\frac{\ii}{2}ab}P_\Phi. \tag{$\dagger$}
  \]
  \tutornote{式 $(\dagger)$ 的指数项底本原文作 $\frac{\ii}{2}ab$，缺少因子 $\hbar$；而此前一行已算得 $\frac{\ii\hbar}{2}ab$。此处照原文录入，未擅改——取 $a=b=0$ 时两种写法都给出 $1$，故不影响本节余下论证。}
  特别地，取 $a=b=0$ 得
  \[
  P_\Phi^2=P_\Phi,\qquad P_\Phi^*=P_\Phi,
  \]
  故 $P_\Phi$ 定义正交投影。可以证明 $P_\Phi\ne0$（请读者自证）。

  \textbf{断言}：$\dim P_\Phi(\mathcal{H})=1$。

  事实上，设有两个向量 $\xi_1,\xi_2\in P_\Phi(\mathcal{H})$。不妨设 $\langle\xi_1|\xi_2\rangle=0$。则对任意 $t_i,s_i\in\R$，式 $(\dagger)$ 给出
  \[
  \begin{aligned}
  \langle U(t_1)V(s_1)\xi_1|U(t_2)V(s_2)\xi_2\rangle
  &=\langle\xi_1|P_\Phi V(-s_1)U(-t_1)U(t_2)V(s_2)P_\Phi\xi_2\rangle\\
  &\propto\langle\xi_1|\xi_2\rangle=0.
  \end{aligned}
  \]
  给定任意向量 $\xi\in P_\Phi(\mathcal{H})$，记
  \[
  \mathcal{H}_\xi:=\Span\{U(t)V(s)\xi\mid t,s\in\R\}.
  \]
  于是 $\mathcal{H}_{\xi_1}\perp\mathcal{H}_{\xi_2}$，且每个 $\mathcal{H}_{\xi_i}$ 给出海森堡群的一个表示。这与 $\mathcal{H}$ 的不可约性矛盾。因此
  \[
  P_\Phi(\mathcal{H})=\Span\{\xi\},\qquad \mathcal{H}=\mathcal{H}_\xi.
  \]
  此外，内积具有形式
  \[
  \begin{aligned}
  \langle U(t_1)V(s_1)\xi|U(t_2)V(s_2)\xi\rangle
  &=\langle\xi|P_\Phi V(-s_1)U(-t_1)U(t_2)V(s_2)P_\Phi\xi\rangle\\
  &=\alpha(t_1,s_1,t_2,s_2)\langle\xi|\xi\rangle,
  \end{aligned}
  \]
  其中 $\alpha(t_1,s_1,t_2,s_2)$ 完全由式 $(\dagger)$ 确定。由此 $\mathcal{H}=\mathcal{H}_\xi$ 必与薛定谔表示酉等价。
\end{proof}

% [原文 p.228]
\begin{example}{4.4.2}
  在薛定谔表示中考虑 $P_\Phi$。对任意函数 $\varphi(x)$，
  \begin{align*}
  (P_\Phi\varphi)(x)
  &=\frac{\hbar}{2\pi}\int_{\R^2}\ee^{-\frac{\ii}{2}\hbar ts}\ee^{-\frac{\hbar}{4}(t^2+s^2)}(U(t)V(s)\varphi)(x)\,\rd t\,\rd s\\
  &=\frac{\hbar}{2\pi}\int_{\R^2}\ee^{-\frac{\ii}{2}\hbar ts}\ee^{-\frac{\hbar}{4}(t^2+s^2)}
  \ee^{\ii\hbar ts}\ee^{\ii sx}\varphi(x+\hbar t)\,\rd t\,\rd s\\
  &=\frac{\hbar}{2\pi}\int_{\R^2}\ee^{-\frac{\hbar}{4}\left(s-\frac{2\ii x}{\hbar}-\ii t\right)^2}
  \ee^{-\frac{\hbar}{2}\left(t+\frac{x}{\hbar}\right)^2}\ee^{-\frac{x^2}{2\hbar}}\varphi(x+\hbar t)\,\rd t\,\rd s\\
  &\propto\ee^{-\frac{x^2}{2\hbar}}.
  \end{align*}
  所以 $P_\Phi$ 是到谐振子基态上的投影。
\end{example}

\takeaway{本节证明斯通–冯·诺依曼定理：满足外尔关系的不可约酉表示唯一，必与薛定谔表示酉等价。证明用外尔关系造出投影 $P_\Phi$，再证其秩为一，从而内积完全确定。}

\section{薛定谔算子}

我们将应用谱理论来研究薛定谔算子
\[
H=H_0+V=-\Delta+V,
\]
其中 $\Delta=\sum_{i=1}^n\frac{\pa}{\pa x^i}\frac{\pa}{\pa x^i}$ 是拉普拉斯算子，$V(\mathbf{x})$ 是势函数。

\subsection{自由粒子}

考虑 $\R^n$ 中运动的自由粒子。哈密顿算子是
\[
H_0=-\Delta.
\]
在傅里叶变换
\[
F:L^2(\R^n)\longrightarrow L^2(\R^n),\qquad
\varphi(\mathbf{x})\longmapsto\widehat{\varphi}(\boldsymbol{\xi})
=\int_{\R^n}\ee^{-\ii\langle\mathbf{x},\boldsymbol{\xi}\rangle}\varphi(\mathbf{x})\,\rd\mathbf{x}
\]
下，算子 $-\Delta$ 变成乘法算子 $M_{\xi^2}$。由于 $M_{\xi^2}$ 自伴、$\xi^2$ 的本质值域是 $[0,\infty)$，可知 $-\Delta$ 是自伴算子，且
\[
\Dom(-\Delta)=\left\{\varphi(\mathbf{x})\in L^2(\R^n)\ \middle|\ \xi^2\widehat{\varphi}(\boldsymbol{\xi})\in L^2(\R^n)\right\},
\]
并且
\[
\sigma(-\Delta)=[0,\infty)=\sigma_{\mathrm{ess}}(-\Delta).
\]
我们用傅里叶变换计算 $\varphi(\mathbf{x})\in L^2(\R^n)$ 的谱测度 $\rd\varphi_{-\Delta,\varphi}(\lambda)$：
\begin{align*}
\langle\varphi|R_{-\Delta}(z)\varphi\rangle
&=\frac{1}{(2\pi)^n}\int_{\R^n}\overline{\widehat{\varphi}(\boldsymbol{\xi})}
\frac{1}{\xi^2-z}\widehat{\varphi}(\boldsymbol{\xi})\,\rd\boldsymbol{\xi}\\
&=\frac{1}{(2\pi)^n}\int\frac{1}{r^2-z}\abs{\widehat{\varphi}(r\Omega)}^2r^{n-1}\,\rd r\,\rd\Omega\\
&=\frac{1}{(2\pi)^n}\int_0^{+\infty}\frac{\lambda^{\frac n2-1}}{2(\lambda-z)}
\left(\int_{S^{n-1}}\abs{\widehat{\varphi}(\sqrt{\lambda}\Omega)}^2\,\rd\Omega\right)\rd\lambda.
\end{align*}

% [原文 p.229]
因此
\[
\rd\varphi_{-\Delta,\varphi}(\lambda)
=\frac{1}{2\cdot(2\pi)^n}\chi_{[0,\infty)}(\lambda)\lambda^{\frac n2-1}
\left(\int_{S^{n-1}}\abs{\widehat{\varphi}(\sqrt{\lambda}\Omega)}^2\,\rd\Omega\right)\rd\lambda.
\]
预解式 $R_{-\Delta}(z)$ 可直接用热算子算出：
\[
R_{-\Delta}(z)=\int_0^\infty\ee^{zt}\ee^{t\Delta}\,\rd t,\qquad \Real z<0.
\]
热算子 $\ee^{t\Delta}$ 的核是
\[
h_t(\mathbf{x},\mathbf{y})=\frac{1}{(4\pi t)^{\frac n2}}\ee^{-\frac{\abs{\mathbf{x}-\mathbf{y}}^2}{4t}}.
\]
于是 $R_{-\Delta}(z)$ 的核是
\[
G_z(\mathbf{x},\mathbf{y})=\int_0^\infty\frac{1}{(4\pi t)^{\frac n2}}
\ee^{-\frac{\abs{\mathbf{x}-\mathbf{y}}^2}{4t}+zt}\,\rd t,\qquad \Real z<0,
\]
它可用修改的贝塞尔函数表达为
\[
G_z(\mathbf{x},\mathbf{y})
=\frac{1}{2\pi}\left(\frac{\sqrt{-z}}{2\pi r}\right)^{\frac n2-1}
K_{\frac n2-1}\left(\sqrt{-z}\,r\right),
\qquad r=\abs{\mathbf{x}-\mathbf{y}},\quad \Real z<0.
\]

\subsection{谐振子}

考虑一维谐振子
\[
A=-\frac12\Delta+\frac12\omega^2x^2
=-\frac12\left(\frac{\rd}{\rd x}\right)^2+\frac12\omega^2x^2.
\]
用阶梯算子
\[
a^\dagger=\frac{1}{\sqrt{2\omega}}\left(\omega x-\frac{\rd}{\rd x}\right),
\qquad
a=\frac{1}{\sqrt{2\omega}}\left(\omega x+\frac{\rd}{\rd x}\right),
\]
可将其表示为
\[
A=\omega\left(a^\dagger a+\frac12\right).
\]
由此得到基态
\[
\psi_0(x)=\left(\frac{\omega}{\pi}\right)^{\frac14}\ee^{-\frac{\omega}{2}x^2},
\]
以及第 $n$ 个激发态
\[
\psi_n(x)=\frac{1}{\sqrt{n!}}(a^\dagger)^n\psi_0(x)
=\left(\frac{\omega}{\pi}\right)^{\frac14}\frac{1}{\sqrt{2^nn!}}H_n(\sqrt{\omega}x)\ee^{-\frac{\omega}{2}x^2}.
\]
这里
\[
H_n(y)=(-1)^n\ee^{y^2}\frac{\rd^n}{\rd y^n}\left(\ee^{-y^2}\right)
\]
是厄米多项式。本征态 $\{\psi_n\}_{n=0}^\infty$ 构成 $L^2(\R)$ 中的正交归一系。我们证明它们其实构成一组基。设 $\varphi\in(\Span\{\psi_n\}_{n=0}^\infty)^\perp$。则
\[
\int_\R x^n\varphi(x)\ee^{-\frac{\omega}{2}x^2}=0,\qquad\forall n=0,1,\cdots.
\]

% [原文 p.230]
注意
\[
\sum_{n=0}^\infty\frac{(-\ii\xi x)^n}{n!}\ee^{-\frac{\omega}{2}x^2}
=\ee^{-\ii\xi x}\ee^{-\frac{\omega}{2}x^2}
\]
在 $L^2(\R)$ 中成立。于是
\[
\int_\R\ee^{-\ii\xi x}\varphi(x)\ee^{-\frac{\omega}{2}x^2}=0,\qquad\forall\xi\in\R.
\]
这说明 $\varphi(x)\ee^{-\frac{\omega}{2}x^2}$ 的傅里叶变换为零，故 $\varphi(x)\ee^{-\frac{\omega}{2}x^2}=0$，即 $\varphi=0\in L^2(\R)$。

由于 $\{\psi_n\}_{n=0}^\infty$ 构成 $L^2(\R)$ 的正交归一基，且
\[
A\psi_n=\left(n+\frac12\right)\omega\psi_n,
\]
我们有
\[
\sigma(A)=\left\{\left(n+\frac12\right)\omega\right\}_{n=0}^\infty=\sigma_{\mathrm{disc}}(A).
\]

\subsection{半有界势}

\begin{definition}{4.5.1}
  设 $T\in\mathcal{D}'(\R^n)$ 是分布。若对任意非负检验函数 $f\ge0$ 都有
  \[
  T(f)\ge0,
  \]
  则称 $T$ 是\emph{正的}，记作 $T\ge0$。对 $T_1,T_2\in\mathcal{D}'(\R^n)$，若 $T_1-T_2\ge0$，则记 $T_1\ge T_2$。
\end{definition}

\begin{example}{4.5.2}
  考虑分布
  \[
  T_\varphi(f)=\int_{\R^n}\varphi(\mathbf{x})f(\mathbf{x})\,\rd\mathbf{x},
  \]
  其中 $\varphi(\mathbf{x})$ 是连续函数。则 $T_\varphi\ge0$ 当且仅当对一切 $\mathbf{x}\in\R^n$ 有 $\varphi(\mathbf{x})\ge0$。
\end{example}

\begin{lemma}{4.5.3}
  若 $f_n\to f$ 于 $L^1$，则存在子列 $f_{n_k}$，它几乎处处逐点收敛到 $f$。
\end{lemma}

\begin{proof}
  选取子列 $f_{n_k}$，使得
  \[
  \norm{f_{n_k}-f_{n_l}}_{L^1}<\frac{1}{2^N},\qquad\forall k,l\ge N.
  \]
  考虑非负函数
  \[
  g_k=\abs{f_{n_1}}+\sum_{m=1}^k\abs{f_{n_{m+1}}-f_{n_m}}.
  \]
  我们有
  \[
  \int g_k(\mathbf{x})\,\rd\mathbf{x}
  =\norm{f_{n_1}}_{L^1}+\sum_{m=1}^k\norm{f_{n_{m+1}}-f_{n_m}}_{L^1}
  \le\norm{f_{n_1}}_{L^1}+2.
  \]
  由单调收敛，$g_k\to g$ 于 $L^1$，故 $g<\infty$ 几乎处处。这说明 $f_{n_k}\to f$ 几乎处处。
\end{proof}

% [原文 p.231]
\begin{definition}{4.5.4}
  设 $p\ge1$。定义\emph{局部 $L^p$ 空间}
  \[
  L^p_{\mathrm{loc}}(\R^n)=\left\{f\ \middle|\ 
  \int_\Omega\abs{f(\mathbf{x})}^p\,\rd\mathbf{x}<\infty,\ 
  \forall\ \text{有界 }\Omega\subset\R^n\right\}.
  \]
\end{definition}
显然 $L^p(\R^n)\subset L^p_{\mathrm{loc}}(\R^n)$。

\begin{theorem}{4.5.5}（加藤不等式）
  设 $u\in L^1_{\mathrm{loc}}(\R^n)$ 且 $\Delta u\in L^1_{\mathrm{loc}}(\R^n)$。则有分布意义下的不等式
  \[
  \Delta\abs{u}\ge\Real\left(\sgn(u)\Delta u\right).
  \]
  这里
  \[
  (\sgn u)(\mathbf{x})=\begin{cases}
  0,&u(\mathbf{x})=0,\\[2mm]
  \dfrac{\overline{u(\mathbf{x})}}{\abs{u(\mathbf{x})}},&u(\mathbf{x})\ne0.
  \end{cases}
  \]
\end{theorem}

\begin{proof}
  注意由假设 $\Delta u\in L^1_{\mathrm{loc}}(\R^n)$，且 $\sgn u$ 有界，故 $\sgn(u)\Delta u\in L^1_{\mathrm{loc}}(\R^n)$，因而也是分布。

  先考虑 $u$ 为 $C^\infty$ 的情形。令
  \[
  \varphi_\epsilon(\mathbf{x})=\sqrt{\abs{u(\mathbf{x})}^2+\epsilon^2},\qquad\epsilon>0.
  \]
  对 $\varphi_\epsilon^2$ 作用 $\nabla$ 与 $\Delta$，得
  \[
  \varphi_\epsilon\nabla\varphi_\epsilon=\Real\left(\bar u\nabla u\right),
  \]
  \[
  \varphi_\epsilon\Delta\varphi_\epsilon+\abs{\nabla\varphi_\epsilon}^2
  =\Real\left(\bar u\Delta u\right)+\abs{\nabla u}^2.
  \]
  由于 $\varphi_\epsilon\ge\abs{u}$，
  \[
  \abs{\nabla\varphi_\epsilon}
  =\abs{\Real\left(\frac{\bar u}{\varphi_\epsilon}\nabla u\right)}
  \le\abs{\nabla u}
  \ \Longrightarrow\
  \varphi_\epsilon\Delta\varphi_\epsilon\ge\Real\left(\bar u\Delta u\right),
  \]
  即
  \[
  \Delta\varphi_\epsilon\ge\Real\left(\frac{\bar u}{\varphi_\epsilon}\Delta u\right)
  \]
  逐点成立。取极限 $\epsilon\to0$，得到分布不等式
  \[
  \Delta\abs{u}\ge\Real\left(\sgn(u)\Delta u\right).
  \]

  现在考虑 $u\in L^1_{\mathrm{loc}}(\R^n)$ 且 $\Delta u\in L^1_{\mathrm{loc}}(\R^n)$ 的情形。取 $\rho\in C_c^\infty(\R^n)$ 使
  \[
  \rho\ge0,\qquad\int_{\R^n}\rho(\mathbf{x})\,\rd\mathbf{x}=1.
  \]
  令 $\rho_\delta(\mathbf{x})=\frac{1}{\delta^n}\rho\left(\frac{\mathbf{x}}{\delta}\right)$，并记
  \[
  u_\delta(\mathbf{x})=(\rho_\delta*u)(\mathbf{x})
  =\int_{\R^n}\rho_\delta(\mathbf{x}-\mathbf{y})u(\mathbf{y})\,\rd\mathbf{y}.
  \]
  则 $u_\delta$ 光滑。令
  \[
  \varphi_{\epsilon,\delta}=\sqrt{\abs{u_\delta}^2+\epsilon^2}.
  \]
% [原文 p.232]
  则如上可得逐点不等式
  \[
  \Delta\varphi_{\epsilon,\delta}\ge\Real\left(\frac{\overline{u_\delta}}{\varphi_{\epsilon,\delta}}\Delta u_\delta\right).
  \]
  由于 $u_\delta\to u$ 于局部 $L^1$ 范数，由命题 4.5.3 与对角线论证，可取子列 $u_{\delta_k}\to u$ 在 $\delta_k\to0$ 时几乎处处逐点收敛。注意 $\Delta u_\delta=(\Delta u)_\delta$，且 $\Delta u\in L^1_{\mathrm{loc}}(\R^n)$，故 $\Delta u_\delta\to\Delta u$ 于 $L^1_{\mathrm{loc}}(\R^n)$，从而也在分布意义下收敛。于是令 $\delta_k\to0$，
  \[
  \Delta\varphi_\epsilon\ge\Real\left(\frac{\bar u}{\varphi_\epsilon}\Delta u\right)
  \]
  在分布意义下成立。再令 $\epsilon\to0$，利用 $\left|\frac{\bar u}{\varphi_\epsilon}\right|\le1$ 一致有界以及 $\Delta u\in L^1_{\mathrm{loc}}(\R^n)$，得到
  \[
  \Delta\varphi\ge\Real\left(\sgn(u)\Delta u\right)
  \]
  在分布意义下成立。
  \tutornote{原文此处写作 $\Delta\varphi$，由上下文 $\varphi_\epsilon\to\abs{u}$，疑应为 $\Delta\abs{u}$。}
\end{proof}

\begin{theorem}{4.5.6}
  设 $V(\mathbf{x})\in L^2_{\mathrm{loc}}(\R^n)$ 且逐点 $V\ge0$。则薛定谔算子
  \[
  H=-\Delta+V\quad\text{其中}\quad\Dom(H)=C_c^\infty(\R^n)
  \]
  在 $L^2(\R^n)$ 上本质自伴。
\end{theorem}

\begin{proof}
  显然 $H$ 是 $C_c^\infty(\R^n)$ 上的对称算子。由命题 4.1.24，只需证明 $\Ker(H^*+1)=\{0\}$。设 $u\in\Dom(H^*)$ 满足 $(H^*+1)u=0$。这等价于 $u\in L^2(\R^n)$ 满足分布方程
  \[
  -\Delta u+Vu+u=0.
  \]
  由施瓦茨不等式，
  \[
  V\in L^2_{\mathrm{loc}},\quad u\in L^2
  \ \Longrightarrow\ Vu\in L^1_{\mathrm{loc}}.
  \]
  利用包含关系 $L^2\subset L^2_{\mathrm{loc}}\subset L^1_{\mathrm{loc}}$，得
  \[
  \Delta u=Vu+u\in L^1_{\mathrm{loc}}.
  \]
  现在可以应用加藤不等式，得到
  \[
  \Delta\abs{u}\ge\Real(\sgn(u)\Delta u)
  =\Real(\sgn(u)(Vu+u))=(V+1)\abs{u}\ge0.
  \]
  取 $\rho\in C_c^\infty(\R^n)$ 使
  \[
  \rho\ge0,\qquad\int_{\R^n}\rho(\mathbf{x})\,\rd\mathbf{x}=1.
  \]
  令 $\rho_\epsilon(\mathbf{x})=\frac{1}{\epsilon^n}\rho\left(\frac{\mathbf{x}}{\epsilon}\right)$，并记
  \[
  \varphi_\epsilon(\mathbf{x})=(\rho_\epsilon*\abs{u})(\mathbf{x})
  =\int_{\R^n}\rho_\epsilon(\mathbf{x}-\mathbf{y})\abs{u(\mathbf{y})}\,\rd\mathbf{y},
  \qquad\epsilon>0,
  \]
% [原文 p.233]
  它是光滑的。由于
  \[
  \Delta\varphi_\epsilon=(\Delta\rho_\epsilon)*\abs{u}\in L^2,
  \]
  我们有 $\varphi_\epsilon\in\Dom(\Delta)$，从而由傅里叶变换知 $\langle\varphi_\epsilon|\Delta\varphi_\epsilon\rangle\le0$。

  另一方面，由上面的 $\Delta\abs{u}\ge0$，有分布不等式
  \[
  \Delta\varphi_\epsilon=\rho_\epsilon*\Delta\abs{u}\ge0.
  \]
  注意 $\varphi_\epsilon\ge0$ 逐点成立。于是 $\langle\varphi_\epsilon|\Delta\varphi_\epsilon\rangle\ge0$，从而
  \[
  \langle\varphi_\epsilon|\Delta\varphi_\epsilon\rangle=0.
  \]
  这说明 $\varphi_\epsilon=0$。于是 $\abs{u}=\lim_{\epsilon\to0^+}\varphi_\epsilon=0$ 于 $L^2$。
\end{proof}

\begin{example}{4.5.7}
  ① 算子
  \[
  -\Delta+\norm{\mathbf{x}}^2,\qquad -\Delta+\norm{\mathbf{x}}^2+\norm{\mathbf{x}}^4
  \]
  本质自伴。
  ② 设 $V(\mathbf{x})=\frac{1}{\norm{\mathbf{x}}}$ 于 $\R^3$。则
  \[
  -\Delta+V=-\Delta+\frac{1}{\norm{\mathbf{x}}}
  \]
  本质自伴。
\end{example}

\subsection{加藤–雷利希定理}

\begin{definition}{4.5.8}
  设 $B$ 是算子。若
  \begin{enumerate}[label=(\arabic*),leftmargin=3.2em,itemsep=1pt]
    \item $\Dom(A)\subset\Dom(B)$；
    \item 存在 $a,b\ge0$ 使得
    \[
    \norm{B\varphi}\le a\norm{A\varphi}+b\norm{\varphi},\qquad\forall\varphi\in\Dom(A),
    \]
  \end{enumerate}
  则称 $B$ \emph{相对于 $A$ 有界}。记作 $B$ 是 $A$-有界的。满足上述不等式的 $a$ 的下确界称为 $B$ 的 \emph{$A$-界}。
\end{definition}

\begin{theorem}{4.5.9}（加藤–雷利希定理）
  设 $A$ 自伴，$B$ 对称，且 $B$ 相对于 $A$ 有界、$A$-界小于 $1$。则
  \begin{enumerate}[label=(\arabic*),leftmargin=3.2em,itemsep=1pt]
    \item $A+B$ 在 $\Dom(A)$ 上自伴；
    \item 设 $\norm{B\varphi}\le a\norm{A\varphi}+b\norm{\varphi}$（$a<1$）且 $A\ge c$。则
    \[
    A+B\ge c-\max\left\{a\abs{c}+b,\ \frac{b}{1-a}\right\}.
    \]
  \end{enumerate}
\end{theorem}

% [原文 p.234]
\begin{proof}
  ① 设 $\lambda>0$。对任意 $\varphi\in\Dom(A)$，
  \[
  \norm{(A+\ii\lambda)\varphi}^2=\norm{A\varphi}^2+\lambda^2\norm{\varphi}^2
  \ \Longrightarrow\
  \norm{A(A+\ii\lambda)^{-1}}\le1,
  \]
  且
  \[
  \norm{(A+\ii\lambda)^{-1}}\le\frac{1}{\lambda}.
  \]
  由 $B$ 的 $A$-有界性，对任意 $\psi\in\mathcal{H}$，
  \[
  \norm{B(A+\ii\lambda)^{-1}\psi}
  \le a\norm{A(A+\ii\lambda)^{-1}\psi}+b\norm{(A+\ii\lambda)^{-1}\psi}
  \le\left(a+\frac{b}{\lambda}\right)\norm{\varphi}.
  \]
  \tutornote{上式中右端的 $\norm{\varphi}$ 与左端的 $\psi$ 不一致，疑原文笔误，应为 $\norm{\psi}$。}
  由于 $a<1$，可取 $\lambda$ 充分大使 $a+\frac{b}{\lambda}<1$。对此 $\lambda$，
  \[
  \norm{B(A+\ii\lambda)^{-1}}<1.
  \]
  故 $\Ran\left(I+B(A+\ii\lambda)^{-1}\right)=\mathcal{H}$。由于
  \[
  B+A+\ii\lambda=\left(I+B(A+\ii\lambda)^{-1}\right)(A+\ii\lambda),
  \]
  且由 $A$ 自伴知 $\Ran(A+\ii\lambda)=\mathcal{H}$，我们有
  \[
  \Ran(B+A+\ii\lambda)=\mathcal{H}.
  \]
  由命题 4.1.20，$B+A$ 自伴。

  ② 设 $t>-c$，则 $-t\notin\sigma(A)$ 且由命题 4.1.24 有 $\Ran(A+t)=\mathcal{H}$。现在对
  \[
  t+c>\max\left\{a\abs{c}+b,\ \frac{b}{1-a}\right\},
  \]
  我们有
  \begin{align*}
  \norm{B(A+t)^{-1}\psi}
  &\le a\norm{A(A+t)^{-1}\psi}+b\norm{(A+t)^{-1}\psi}\\
  &\le\left(a\max\left\{\frac{\abs{c}}{c+t},1\right\}+\frac{b}{c+t}\right)\norm{\psi}\\
  &<\norm{\psi}
  \end{align*}
  从而
  \[
  \norm{B(A+t)^{-1}}<1.
  \]
  因此对这种 $t$，
  \[
  (A+B+t)^{-1}=(A+t)^{-1}\left(1+B(A+t)^{-1}\right)^{-1},
  \]
  于是 $-t\in\rho(A+B)$。
\end{proof}

\begin{proposition}{4.5.10}
  设 $V=V_1+V_2$ 是实值函数，其中 $V_1\in L^2(\R^3)$，$V_2\in L^\infty(\R^3)$。则 $-\Delta+V$ 在 $D(-\Delta)$ 上自伴。
\end{proposition}

\begin{proof}
  令 $A=-\Delta$，$B=V=V_1+V_2$。对任意 $\varphi\in C_c^\infty(\R^3)$，设 $\widehat{\varphi}$ 是其傅里叶变换
  \[
  \varphi(\mathbf{x})=\frac{1}{(2\pi)^3}\int_{\R^3}
  \ee^{\ii\langle\mathbf{x},\boldsymbol{\xi}\rangle}\widehat{\varphi}(\boldsymbol{\xi})\,\rd\boldsymbol{\xi}.
  \]
% [原文 p.235]
  对任意 $\mathbf{x}\in\R^3$，
  \begin{align*}
  \abs{\varphi(\mathbf{x})}
  &\le\frac{1}{(2\pi)^3}\int_{\R^3}
  \frac{(1+\lambda^2\xi^2)\abs{\widehat{\varphi}(\boldsymbol{\xi})}}{(1+\lambda^2\xi^2)}\,\rd\boldsymbol{\xi}
  \qquad\lambda>0\\
  &\le\frac{1}{(2\pi)^3}
  \left(\int_{\R^3}(1+\lambda^2\xi^2)^2\abs{\widehat{\varphi}(\boldsymbol{\xi})}^2\,\rd\boldsymbol{\xi}\right)^{\frac12}
  \left(\int_{\R^3}\frac{1}{(1+\lambda^2\xi^2)^2}\,\rd\boldsymbol{\xi}\right)^{\frac12}\\
  &=\frac{c}{\lambda^{3/2}}
  \left(\int_{\R^3}(1+\lambda^2\xi^2)^2\abs{\widehat{\varphi}(\boldsymbol{\xi})}^2\,\rd\boldsymbol{\xi}\right)^{\frac12}\\
  &\le a\norm{\Delta\varphi}_{L^2}+b\norm{\varphi}_{L^2},
  \end{align*}
  其中 $a$ 可通过调节 $\lambda$ 取得任意小。因此
  \begin{align*}
  \norm{V\varphi}_{L^2}
  &\le\norm{V_1\varphi}_{L^2}+\norm{V_2\varphi}_{L^2}\\
  &\le\norm{V_1}_{L^2}\norm{\varphi}_{L^\infty}+\norm{V_2}_{L^\infty}\norm{\varphi}_{L^2}\\
  &\le a\norm{V_1}_{L^2}\norm{\Delta\varphi}_{L^2}
  +\left(b+\norm{V_2}_{L^\infty}\right)\norm{\varphi}_{L^2}.
  \end{align*}
  这说明 $\Dom(V)\supset\Dom(-\Delta)$。可取 $a$ 充分小使 $a\norm{V_1}_{L^2}<1$。由加藤–雷利希定理，$-\Delta+V$ 自伴。
\end{proof}

\begin{example}{4.5.11}
  考虑 $\R^3$ 上的下述库仑势
  \[
  V(\mathbf{x})=-\frac{e^2}{\norm{\mathbf{x}}}.
  \]
  则 $-\Delta+V$ 自伴。事实上，令
  \[
  V_1(\mathbf{x})=\begin{cases}
  V(\mathbf{x}),&\norm{\mathbf{x}}\le1,\\
  0,&\norm{\mathbf{x}}>1,
  \end{cases}
  \qquad V_2(\mathbf{x})=V(\mathbf{x})-V_1(\mathbf{x}).
  \]
  则 $V_1\in L^2(\R^3)$，$V_2\in L^\infty(\R^3)$，加藤–雷利希定理适用。
\end{example}

\subsection{特罗特乘积公式与费曼积分}

我们经常需要处理形如 $A+B$ 的算子。例如哈密顿量
\[
H=H_0+V,
\]
其中 $H_0=-\Delta$ 是自由部分，$V$ 编码势函数。

\begin{proposition}{4.5.12}（特罗特乘积公式）
  设 $A,B$ 是自伴算子，且 $A+B$ 自伴。则
  \[
  \ee^{\ii t(A+B)}=\slim_{n\to\infty}
  \left[\ee^{\ii tA/n}\ee^{\ii tB/n}\right]^n.
  \]
\end{proposition}

\begin{proof}
  这是李乘积公式的无穷维版本。把恒等式
  \[
  \begin{aligned}
  X^n-Y^n&=X^n-X^{n-1}Y+X^{n-1}Y-X^{n-2}Y^2+\cdots+XY^{n-1}-Y^n\\
  &=\sum_{k=0}^{n-1}X^k(X-Y)Y^{n-1-k}.
  \end{aligned}
  \]
% [原文 p.236]
  用于 $X=\ee^{\frac{\ii}{n}tA}\ee^{\frac{\ii}{n}tB}$、$Y=\ee^{\frac{\ii}{n}t(A+B)}$ 以及 $\varphi\in\mathcal{H}$，得
  \begin{align*}
  &\left\|\left[\left(\ee^{\frac{\ii}{n}tA}\ee^{\frac{\ii}{n}tB}\right)^n-\ee^{\ii t(A+B)}\right]\varphi\right\|\\
  &\le n\max_{0\le k\le n-1}
  \left\|\left(\ee^{\frac{\ii}{n}tA}\ee^{\frac{\ii}{n}tB}-\ee^{\frac{\ii}{n}t(A+B)}\right)
  \ee^{\frac{\ii}{n}tk(A+B)}\varphi\right\|\\
  &=t\max_{0\le k\le n-1}
  \left\|G\left(\frac tn\right)\ee^{\frac{\ii}{n}tk(A+B)}\varphi\right\|\\
  &\le t\max_{-1\le\tau\le1}
  \left\|G\left(\frac tn\right)\ee^{\ii\tau(A+B)}\varphi\right\|,
  \end{align*}
  其中 $G(s)=\frac1s\left(\ee^{\ii sA}\ee^{\ii sB}-\ee^{\ii s(A+B)}\right)$。

  由于 $A+B$ 自伴，$\Dom(A+B)$ 在图像范数
  \[
  \norm{\varphi}_{\Gamma(A+B)}=\norm{\varphi}+\norm{(A+B)\varphi}
  \]
  下是希尔伯特空间。对 $\varphi\in\Dom(A+B)$，
  \begin{align*}
  \lim_{s\to0}G(s)\cdot\varphi
  &=\lim_{s\to0}\left(\ee^{\ii sA}\left(\frac{\ee^{\ii sB}-I}{s}\right)
  +\frac{\ee^{\ii sA}-I}{s}+\frac{I-\ee^{\ii s(A+B)}}{s}\right)\varphi\\
  &=(\ii B+\ii A-\ii(A+B))\varphi=0,\\
  \lim_{s\to\infty}G(s)\cdot\varphi&=0.
  \end{align*}
  所以 $\norm{G(s)\varphi}$ 对 $s\in\R$ 有界。由一致有界性定理，存在常数 $C>0$ 使得
  \[
  \norm{G(s)\varphi}\le C\norm{\varphi}_{\Gamma(A+B)},\qquad
  \forall\varphi\in\Dom(A+B).
  \]
  又映射 $\tau\mapsto\ee^{\ii\tau(A+B)}\varphi$ 在 $\norm{\cdot}_{\Gamma(A+B)}$-拓扑下连续，因此族 $\{\ee^{\ii\tau(A+B)}\varphi\}_{-1\le\tau\le1}$ 在 $\Dom(A+B)$ 中紧。由此
  \[
  \lim_{s\to0}\max_{-1\le\tau\le1}G(s)\ee^{\ii\tau(A+B)}\varphi=0
  \]
  对 $\varphi\in\Dom(A+B)$ 成立。命题得证，因为 $\Dom(A+B)$ 稠密。
\end{proof}

特罗特乘积公式与量子力学中的费曼路径积分密切相关。考虑 $\R^n$ 中的量子粒子，且
\[
H=-\frac12\Delta+V=H_0+V,\qquad
\Delta=\sum_{i=1}^n\frac{\pa}{\pa x^i}\frac{\pa}{\pa x^i}.
\]
酉算子 $\ee^{-\ii tH_0}=\ee^{\frac{\ii}{2}t\Delta}$ 是量子微分的指数：
\[
\left(\ee^{\frac{\ii}{2}t\Delta}\varphi\right)(\mathbf{x})
=\frac{\ee^{-\frac{\ii}{4}\pi n}}{(2\pi t)^{n/2}}
\int_{\R^n}\ee^{\frac{\ii\abs{\mathbf{y}}^2}{2t}}\varphi(\mathbf{x}+\mathbf{y})\,\rd\mathbf{y}
=\frac{1}{(2\pi\ii t)^{n/2}}
\int_{\R^n}\ee^{\frac{\ii\abs{\mathbf{x}-\mathbf{y}}^2}{2t}}\varphi(\mathbf{y})\,\rd\mathbf{y}.
\]
若对适当的势 $V(\mathbf{x})$ 应用特罗特乘积公式，我们得到
\begin{align*}
\left(\ee^{-\ii tH}\varphi\right)(\mathbf{x}_0)
&=\lim_{N\to\infty}\left(\ee^{-\ii\frac tN V}\ee^{-\ii\frac tN H_0}\right)^N\varphi(\mathbf{x}_0)\\
&=\lim_{N\to\infty}\left(\frac{N}{2\pi\ii t}\right)^{\frac{Nn}{2}}
\int_{\R^n}\prod_{j=1}^N\rd\mathbf{x}_j\,
\ee^{\ii S_N(\mathbf{x}_0,\cdots,\mathbf{x}_N;t)}\varphi(\mathbf{x}_N),
\end{align*}
% [原文 p.237]
其中
\begin{align*}
S_N(\mathbf{x}_0,\cdots,\mathbf{x}_N;t)
&=\sum_{j=0}^{N-1}\frac{\abs{\mathbf{x}_{j+1}-\mathbf{x}_j}^2}{2t/N}
-\frac tN V(\mathbf{x}_j)\\
&=\frac tN\sum_{j=0}^{N-1}\left[
\frac12\left|\frac{\mathbf{x}_{j+1}-\mathbf{x}_j}{t/N}\right|^2-V(\mathbf{x}_j)\right].
\end{align*}
于是极限 $N\to\infty$ 恰好模拟了 §2.2.4 讨论的费曼路径积分。

\takeaway{本节用谱理论研究薛定谔算子：自由粒子经傅里叶变换化为乘法算子，谐振子用阶梯算子求得离散谱；再用加藤不等式与加藤–雷利希定理给出势的自伴性判据，最后用特罗特公式导出费曼积分。}

\section{弗雷德霍姆算子与指标}

\subsection{弗雷德霍姆算子}

\begin{definition}{4.6.1}
  设 $A$ 是 $\mathcal{H}$ 上的有界算子。若 $\Ker A$ 与 $\operatorname{Coker}A=\mathcal{H}/\Ran A$ 都是有限维的，则称 $A$ 是\emph{弗雷德霍姆算子}。整数
  \[
  \operatorname{ind}(A)=\dim\Ker A-\dim\operatorname{Coker}A
  \]
  称为 $A$ 的\emph{指标}。记
  \[
  \mathbf{F}(\mathcal{H})=\{\mathcal{H}\ \text{上的弗雷德霍姆算子}\}.
  \]
\end{definition}

\begin{example}{4.6.2}
  $\mathcal{H}=\ell^2(\N)=\left\{(a_1,a_2,\cdots)\ \middle|\ 
  a_i\in\C:\ \sum_{i=1}^\infty\abs{a_i}^2<\infty\right\}$。考虑\emph{移位算子}
  \[
  S:(a_1,a_2,\cdots)\longmapsto(0,a_1,a_2,\cdots).
  \]
  则 $\Ker S=0$，$\operatorname{Coker}S=\C$。$S$ 是弗雷德霍姆算子，且 $\operatorname{ind}S=-1$。其伴随
  \[
  S^*:(a_1,a_2,\cdots)\longmapsto(a_2,a_3,\cdots)
  \]
  也是弗雷德霍姆算子，且 $\operatorname{ind}S^*=1$。
\end{example}

\begin{proposition}{4.6.3}
  设 $A,B\in\mathbf{F}(\mathcal{H})$ 是两个弗雷德霍姆算子。则复合 $AB$ 也是弗雷德霍姆算子，且
  \[
  \operatorname{ind}(AB)=\operatorname{ind}(A)+\operatorname{ind}(B).
  \]
\end{proposition}

\begin{proof}
  考虑下述行正合的交换图
  \begin{center}
  \begin{tikzpicture}[>=Stealth,scale=0.86,transform shape,
    every node/.style={inner sep=2pt,font=\small}]
    \node (t0) at (0,0) {$0$};
    \node (t1) at (1.55,0) {$\mathcal{H}$};
    \node (t2) at (4.25,0) {$\mathcal{H}\oplus\mathcal{H}$};
    \node (t3) at (6.95,0) {$\mathcal{H}$};
    \node (t4) at (8.35,0) {$0$};
    \node (b0) at (0,-1.75) {$0$};
    \node (b1) at (1.55,-1.75) {$\mathcal{H}$};
    \node (b2) at (4.25,-1.75) {$\mathcal{H}\oplus\mathcal{H}$};
    \node (b3) at (6.95,-1.75) {$\mathcal{H}$};
    \node (b4) at (8.35,-1.75) {$0$};
    \draw[->] (t0)--(t1);
    \draw[->] (t1)--node[above]{$I\oplus B$}(t2);
    \draw[->] (t2)--node[above]{$(B,-I)$}(t3);
    \draw[->] (t3)--(t4);
    \draw[->] (b0)--(b1);
    \draw[->] (b1)--node[below]{$A\oplus I$}(b2);
    \draw[->] (b2)--node[below]{$(I,-A)$}(b3);
    \draw[->] (b3)--(b4);
    \draw[->] (t1)--node[right]{$B$}(b1);
    \draw[->] (t2)--node[right]{$AB\oplus I$}(b2);
    \draw[->] (t3)--node[right]{$A$}(b3);
  \end{tikzpicture}
  \end{center}
  % [图：原书 p.237 的交换图，本图按原文重绘；图内只有数学符号，原图无题注]
  由五引理即得。
\end{proof}

\begin{proposition}{4.6.4}
  设 $A\in\mathbf{F}(\mathcal{H})$ 是弗雷德霍姆算子。则其伴随 $A^*$ 也是弗雷德霍姆算子，且
  \[
  \operatorname{ind}(A^*)=-\operatorname{ind}A.
  \]
\end{proposition}

\begin{proof}
  $\Ker A^*=(\Ran A)^\perp\simeq\operatorname{Coker}A$，$\Ker A\simeq\operatorname{Coker}A^*$。
\end{proof}

% [原文 p.238]
\begin{theorem}{4.6.5}（阿特金森定理）
  设 $A$ 是 $\mathcal{H}$ 上的有界算子。$A$ 是弗雷德霍姆算子，当且仅当存在有界算子 $B$，使得 $AB-I$ 与 $BA-I$ 都是紧算子。这样的 $B$ 称为 $A$ 的\emph{拟基本解}。
\end{theorem}

\begin{proof}
  设 $A$ 是弗雷德霍姆算子。则
  \[
  \mathcal{H}=\Ker A\oplus(\Ker A)^\perp=\Ran A\oplus\operatorname{Coker}A,
  \]
  且 $A:(\Ker A)^\perp\to\Ran A$ 是有界双射。定义 $B:\Ran A\to(\Ker A)^\perp$ 为其逆，由开映射定理它也有界。再把 $B$ 延拓为 $\mathcal{H}$ 上的有界算子，令 $B|_{\operatorname{Coker}A}=0$。于是
  \[
  I-AB=\ \text{到}\ \operatorname{Coker}A\ \text{上的投影},\qquad
  I-BA=\ \text{到}\ \Ker A\ \text{上的投影}
  \]
  都是有限秩算子，从而是紧算子。

  反之，设
  \[
  K_1=I-AB,\qquad K_2=I-BA
  \]
  都是紧算子。由于
  \[
  K_2=I\quad\text{在}\ \Ker A\ \text{上},
  \]
  由 $K_2$ 的紧性，$\Ker A$ 中的单位球相对紧，故 $\dim\Ker A<\infty$。类似地，利用 $K_1^*=I-B^*A^*$，有
  \[
  \dim\Ker A^*=\dim(\Ran A)^\perp<\infty.
  \]
  只需再证 $\Ran A$ 闭。设 $\varphi_n$ 是 $(\Ker A)^\perp$ 中的序列，且 $A\varphi_n\to\tilde\varphi$。要证 $\tilde\varphi\in\Ran A$。若 $\{\varphi_n\}$ 有界，则由
  \[
  \varphi_n=K_2\varphi_n+BA\varphi_n
  \]
  以及 $K_2$ 紧，可取子列 $\varphi_{n_k}\to\varphi$ 于 $\mathcal{H}$。于是 $\tilde\varphi=\lim_k A\varphi_{n_k}=A\varphi$。

  若 $\{\varphi_n\}$ 无界，则取子列可设 $\norm{\varphi_n}\to\infty$。此时 $\frac{\varphi_n}{\norm{\varphi_n}}$ 有界，且 $A\left(\frac{\varphi_n}{\norm{\varphi_n}}\right)\to0$。重复上面的论证，可得子列
  \[
  \frac{\varphi_n}{\norm{\varphi_n}}\longrightarrow\varphi,\qquad\norm{\varphi}=1,
  \]
  使得 $A\varphi=0$。这与 $\frac{\varphi_n}{\norm{\varphi_n}}\in(\Ker A)^\perp$ 矛盾。
\end{proof}

阿特金森定理说明，弗雷德霍姆算子在紧算子模下可逆。我们可以把这一点说得更精确。

\begin{proposition}{4.6.6}
  若 $A$ 是弗雷德霍姆算子、$K$ 是紧算子，则 $A+K$ 是弗雷德霍姆算子。
\end{proposition}

% [原文 p.239]
\begin{proof}
  取有界算子 $B$ 使 $AB-I$ 与 $BA-I$ 都是紧算子。则
  \[
  (A+K)B-I=(AB-I)+KB,\qquad
  B(A+K)-I=(BA-I)+BK
  \]
  都是紧算子。
\end{proof}

由于 $\mathbf{K}(\mathcal{H})$ 是 $\mathbf{B}(\mathcal{H})$ 中的闭双边理想，商
\[
\mathbf{B}(\mathcal{H})/\mathbf{K}(\mathcal{H})
\]
是巴拿赫代数，称为\emph{卡尔金代数}。弗雷德霍姆算子正是 $\mathbf{B}(\mathcal{H})$ 中映到 $\mathbf{B}(\mathcal{H})/\mathbf{K}(\mathcal{H})$ 的可逆元的那批算子。

\subsection{指标的同伦不变性}

\begin{proposition}{4.6.7}
  $\mathbf{F}(\mathcal{H})$ 是 $\mathbf{B}(\mathcal{H})$ 中的开集，且指标函数
  \[
  \operatorname{ind}:\mathbf{F}(\mathcal{H})\longrightarrow\Z
  \]
  局部常值。
\end{proposition}

\begin{proof}
  可逆性是开条件。设 $A\in\mathbf{F}(\mathcal{H})$。由于
  \[
  A:(\Ker A)^\perp\longrightarrow\Ran A
  \]
  是线性同构，我们用一个希尔伯特空间 $V$ 把它们等同，并写
  \[
  \mathcal{H}=\Ker A\oplus V=V\oplus\operatorname{Coker}A.
  \]
  设 $p:\mathcal{H}=\Ker A\oplus V\to V$ 是投影，$j:V\to\mathcal{H}=V\oplus\operatorname{Coker}A$ 是包含映射。则
  \[
  pAj:V\longrightarrow V
  \]
  是同构。因此对充分接近 $A$ 的 $\widetilde{A}$，
  \[
  p\widetilde{A}j:V\longrightarrow V
  \]
  也是同构。注意 $p$ 与 $j$ 都是弗雷德霍姆算子，且
  \[
  \operatorname{ind}(pAj)=0=\operatorname{ind}(p\widetilde{A}j).
  \]
  由命题 4.6.3，得 $\operatorname{ind}(\widetilde{A})=\operatorname{ind}(A)$。
\end{proof}

\begin{corollary}{4.6.8}
  设 $A$ 弗雷德霍姆、$K$ 紧。则
  \[
  \operatorname{ind}(A+K)=\operatorname{ind}(A).
  \]
\end{corollary}

\begin{proof}
  考虑同伦 $A+tK$。
\end{proof}

% [原文 p.240]
\subsection{托普利茨指标定理}

考虑希尔伯特空间
\[
\mathcal{H}=L^2(S^1),
\]
其中我们等同 $S^1=\{z\in\C\mid\abs{z}=1\}$。内积是
\[
\langle f|g\rangle=\frac{1}{2\pi}\int_{S^1}\overline{f(\theta)}g(\theta)\,\rd\theta.
\]
函数 $z^n$（即 $z^n|_{S^1}=\ee^{\ii n\theta}$）构成 $L^2(S^1)$ 的正交归一基。任意 $f\in L^2(S^1)$ 可展开为
\[
f=\sum_{n\in\Z}\widehat{f}(n)z^n,
\]
其中 $\widehat{f}(n)$ 是第 $n$ 个傅里叶系数
\[
\widehat{f}(n)=\langle z^n|f\rangle
=\frac{1}{2\pi}\int_{S^1}f(\theta)\ee^{-\ii n\theta}\,\rd\theta.
\]
设 $H^2\subset\mathcal{H}$ 是由 $\{z^n\}_{n\ge0}$ 张成的子空间。令
\[
\pi_+:\mathcal{H}\longrightarrow H^2,\qquad
\sum_{n\in\Z}a_nz^n\longmapsto\sum_{n\ge0}a_nz^n
\]
为投影映射。对任意 $f\in C(S^1)$（即 $S^1$ 上的连续函数），定义
\[
T_f:H^2\longrightarrow H^2,\qquad
\varphi\longmapsto\pi_+(f\varphi).
\]
$H^2$ 称为\emph{哈代空间}，$T_f$ 称为与 $f$ 相伴的\emph{托普利茨算子}。

\begin{lemma}{4.6.9}
  对任意 $f,g\in C(S^1)$，$T_f\circ T_g-T_{fg}$ 是紧算子。
\end{lemma}

\begin{proof}
  设 $M_f$ 是 $\mathcal{H}$ 上的乘法算子。注意对任意 $z^n$，交换子
  \[
  [\pi_+,M_{z^n}]
  \]
  是有限秩算子。由于 $f$ 可由 $\sum_{n=-N}^N\widehat{f}(n)z^n$ 逼近，可知
  \[
  [\pi_+,M_f]
  \]
  是紧算子。于是
  \[
  \pi_+M_f\pi_+M_g
  =\pi_+\pi_+M_fM_g+\pi_+[M_f,\pi_+]M_g
  =\pi_+M_{fg}+\text{紧算子},
  \]
  从而
  \[
  T_f\circ T_g=T_{fg}+\text{紧算子}.
  \]
\end{proof}

\begin{corollary}{4.6.10}
  若 $f\in C(S^1)$ 处处不为零，则 $T_f$ 是弗雷德霍姆算子。
\end{corollary}

% [本块止于原书 p.240，下接 ch04c（§4.6.3 的托普利茨指标定理证明与 §4.7）]
